\documentclass[10pt,a4paper]{amsart}
\usepackage{amsmath,amssymb,amsthm,graphicx}
% For commutative diagrams we use Xy-pic
\usepackage[all]{xy}

% We use 2cell for 2-commutative diagrams.
\xyoption{2cell}
% Theorem environments.
%
\theoremstyle{plain}
\newtheorem{theorem}[subsection]{Theorem}
\newtheorem{proposition}[subsection]{Proposition}
\newtheorem{lemma}[subsection]{Lemma}

\theoremstyle{definition}
\newtheorem{definition}[subsection]{Definition}
\newtheorem{example}[subsection]{Example}
\newtheorem{exercise}[subsection]{Exercise}
\newtheorem{situation}[subsection]{Situation}

\theoremstyle{remark}
\newtheorem{remark}[subsection]{Remark}
\newtheorem{remarks}[subsection]{Remarks}

\numberwithin{equation}{subsection}

% Macros
%
\def\lim{\mathop{\mathrm{lim}}\nolimits}
\def\colim{\mathop{\mathrm{colim}}\nolimits}
\def\Spec{\mathop{\mathrm{Spec}}}
\def\Hom{\mathop{\mathrm{Hom}}\nolimits}
\def\Ext{\mathop{\mathrm{Ext}}\nolimits}
\def\SheafHom{\mathop{\mathcal{H}\!\mathit{om}}\nolimits}
\def\SheafExt{\mathop{\mathcal{E}\!\mathit{xt}}\nolimits}
\def\Sch{\mathit{Sch}}
\def\Mor{\mathop{\mathrm{Mor}}\nolimits}
\def\Ob{\mathop{\mathrm{Ob}}\nolimits}
\def\Sh{\mathop{\mathit{Sh}}\nolimits}
\def\NL{\mathop{N\!L}\nolimits}
\def\CH{\mathop{\mathrm{CH}}\nolimits}
\def\proetale{{pro\text{-}\acute{e}tale}}
\def\etale{{\acute{e}tale}}
\def\QCoh{\mathit{QCoh}}
\def\Ker{\mathop{\mathrm{Ker}}}
\def\Im{\mathop{\mathrm{Im}}}
\def\Coker{\mathop{\mathrm{Coker}}}
\def\Coim{\mathop{\mathrm{Coim}}}

% Boxtimes
%
\DeclareMathSymbol{\boxtimes}{\mathbin}{AMSa}{"02}

%
% Macros for moduli stacks/spaces
%
\def\QCohstack{\mathcal{QC}\!\mathit{oh}}
\def\Cohstack{\mathcal{C}\!\mathit{oh}}
\def\Spacesstack{\mathcal{S}\!\mathit{paces}}
\def\Quotfunctor{\mathrm{Quot}}
\def\Hilbfunctor{\mathrm{Hilb}}
\def\Curvesstack{\mathcal{C}\!\mathit{urves}}
\def\Polarizedstack{\mathcal{P}\!\mathit{olarized}}
\def\Complexesstack{\mathcal{C}\!\mathit{omplexes}}
% \Pic is the operator that assigns to X its picard group, usage \Pic(X)
% \Picardstack_{X/B} denotes the Picard stack of X over B
% \Picardfunctor_{X/B} denotes the Picard functor of X over B
\def\Pic{\mathop{\mathrm{Pic}}\nolimits}
\def\Picardstack{\mathcal{P}\!\mathit{ic}}
\def\Picardfunctor{\mathrm{Pic}}
\def\Deformationcategory{\mathcal{D}\!\mathit{ef}}
\usepackage{cleveref}
\crefname{originallemma}{original lemma}{original lemmas}
\Crefname{originallemma}{Original lemma}{Original lemmas}
\makeatletter
\newlabel{topologies-lemma-zariski-etale-smooth-syntomic-fppf}{{021N}{source}{}{sourcepointer-1}{}}
\newlabel{topologies-lemma-zariski-etale-smooth-syntomic-fppf@cref}{{[originallemma][0][]021N}{[1][0][]source}{}{}{}}
\makeatother
\crefname{originallemma}{original lemma}{original lemmas}
\Crefname{originallemma}{Original lemma}{Original lemmas}
\makeatletter
\newlabel{constructions-section-spec-via-glueing}{{01LL}{source}{}{sourcepointer-2}{}}
\newlabel{constructions-section-spec-via-glueing@cref}{{[originallemma][0][]01LL}{[1][0][]source}{}{}{}}
\makeatother
\crefname{originallemma}{original lemma}{original lemmas}
\Crefname{originallemma}{Original lemma}{Original lemmas}
\makeatletter
\newlabel{constructions-section-spec}{{01LQ}{source}{}{sourcepointer-3}{}}
\newlabel{constructions-section-spec@cref}{{[originallemma][0][]01LQ}{[1][0][]source}{}{}{}}
\makeatother
\title{The Kummer sequence is exact for the fppf and syntomic topologies}
\author{The Stacks Project Authors}
\date{stacks-project:etale-cohomology.tex:local-snapshot-20261011; GNU FDL 1.2 or later}
\begin{document}
\maketitle
\noindent\textit{Editorial note.} All delivered mathematical text is the author's own, copied from the pinned Stacks Project chapter file etale-cohomology.tex; no statement or proof step was rewritten, reordered or removed.\par\medskip
\noindent\textit{Editorial note.} The proof is the author's complete argument for both topologies, with the covering argument written out.\par\medskip
\noindent\textit{Editorial note.} Every labelled reference occurring inside the counted spans is carried in assembly.external\_references as a non-load-bearing navigation pointer (same-chapter ones at their definition site, cross-chapter ones at the invocation inside the counted span), so no load-bearing mathematics is replaced by a stub. No bibliography citation occurs inside the counted spans.\par\medskip
\noindent\textit{Editorial note.} The checked chapter list for same-name collisions is the set of stacks chapters staged in this repository: chow, cohomology, divisors, etale-cohomology, flat, more-algebra, more-morphisms, morphisms, spaces-more-morphisms, varieties together with this card's chapter etale-cohomology.\par\medskip
\noindent\textit{Editorial note.} The preamble is the chapter's own native preamble excerpt already used for this source: the Xy-pic setup (native preamble lines 14-18) and the native theorem declarations with the mathematical macros (native preamble lines 154-217).\par\medskip
\noindent\textit{Allowed standard prerequisites (editorial).}\begin{itemize}
\item Sheaf mu\_n of n-th roots of unity and its description as the kernel of raising to the n-th power. Same chapter, the sheaf mu\_n on the big site of S, defined as the kernel of the map (.)\textasciicircum{}n on the multiplicative group. Supplies the first two terms of the sequence and its exactness there by definition.
\item Flat, syntomic and fppf coverings and the local nature of surjectivity of sheaves. Same chapter, the description of the big sites involved and the criterion for a map of sheaves to be surjective by lifting sections locally. Supplies the surjectivity argument, in which a section is lifted after a covering.
\end{itemize}\medskip
\ifdefined\hypertarget\hypertarget{sourcepointer-1}{}\fi
\noindent\textit{Source reference.} Original lemma 021N: sources/stacks-etale-cohomology/pinned-original-source.tex, source lines 2937--2937. Navigation pointer invoked at proof line 2937; the counted span contains the pointer only and the cited result is not reproduced inside it, so the pointer is registered as a non-load-bearing navigation stub in the established form for Stacks material. No bibliography citation is involved.\par\medskip
\ifdefined\hypertarget\hypertarget{sourcepointer-2}{}\fi
\noindent\textit{Source reference.} Original lemma 01LL: sources/stacks-etale-cohomology/pinned-original-source.tex, source lines 2869--2869. Navigation pointer invoked at proof line 2869; the counted span contains the pointer only and the cited result is not reproduced inside it, so the pointer is registered as a non-load-bearing navigation stub in the established form for Stacks material. No bibliography citation is involved.\par\medskip
\ifdefined\hypertarget\hypertarget{sourcepointer-3}{}\fi
\noindent\textit{Source reference.} Original lemma 01LQ: sources/stacks-etale-cohomology/pinned-original-source.tex, source lines 2870--2870. Navigation pointer invoked at proof line 2870; the counted span contains the pointer only and the cited result is not reproduced inside it, so the pointer is registered as a non-load-bearing navigation stub in the established form for Stacks material. No bibliography citation is involved.\par\medskip
\begin{lemma}
\label{lemma-kummer-sequence-syntomic}
For any $n \in \mathbf{N}$ the sequence
$$
0 \to
\mu_{n, S} \to
\mathbf{G}_{m, S} \xrightarrow{(\cdot)^n}
\mathbf{G}_{m, S} \to 0
$$
is a short exact sequence of sheaves on the site
$(\Sch/S)_{fppf}$ and $(\Sch/S)_{syntomic}$.
\end{lemma}

\begin{proof}
By definition the sheaf $\mu_{n, S}$ is the kernel of the map
$(\cdot)^n$. Hence it suffices to show that the last map is surjective.
Since the syntomic topology is weaker than the fppf topology, see
Topologies, Lemma \ref{topologies-lemma-zariski-etale-smooth-syntomic-fppf},
it suffices to prove this for the syntomic topology.
Let $U$ be a scheme over $S$. Let
$f \in \mathbf{G}_m(U) = \Gamma(U, \mathcal{O}_U^*)$.
We need to show that we can find a syntomic cover of
$U$ over the members of which the restriction of $f$ is an $n$th power.
Set
$$
U' =
\underline{\Spec}_U(\mathcal{O}_U[T]/(T^n-f))
\xrightarrow{\pi}
U.
$$
(See
Constructions, Section \ref{constructions-section-spec-via-glueing} or
\ref{constructions-section-spec}
for a discussion of the relative spectrum.)
Let $\Spec(A) \subset U$ be an affine open, and say $f|_{\Spec(A)}$ corresponds
to the unit $a \in A^*$. Then $\pi^{-1}(\Spec(A)) = \Spec(B)$ with
$B = A[T]/(T^n - a)$. The ring map $A \to B$ is finite free of rank $n$,
hence it is faithfully flat, and hence we conclude that
$\Spec(B) \to \Spec(A)$ is surjective. Since this holds for every
affine open in $U$ we conclude that $\pi$ is surjective.
In addition, $B$ is a global relative complete intersection over $A$, so
the ring map $A \to B$ is standard syntomic,
in particular syntomic. Since this holds for every affine open of $U$
we conclude that $\pi$ is syntomic. Hence
$\mathcal{U} = \{\pi : U' \to U\}$ is a syntomic covering.
Moreover, $f|_{U'} = (f')^n$ where $f'$ is the class of $T$
in $\Gamma(U', \mathcal{O}_{U'}^*)$, so $\mathcal{U}$ has the desired property.
\end{proof}

\end{document}
